% Generated draft: requires mathematical and source review.
\documentclass[11pt]{article}
\usepackage[margin=25mm]{geometry}
\usepackage{fontspec}
\setmainfont{Latin Modern Roman}
\usepackage{amsmath,amssymb,mathtools,unicode-math}
\setmathfont{Latin Modern Math}
\usepackage{enumitem,needspace,xurl,hyperref}
\setlength{\parindent}{0pt}
\setlength{\parskip}{5pt}
\setlength{\emergencystretch}{3em}
\newcommand{\N}{\mathbb{N}}
\newcommand{\Z}{\mathbb{Z}}
\newcommand{\Q}{\mathbb{Q}}
\newcommand{\R}{\mathbb{R}}
\newcommand{\C}{\mathbb{C}}
\newcommand{\F}{\mathbb{F}}
\newcommand{\A}{\mathbb{A}}
\newcommand{\PP}{\mathbb{P}}
\newcommand{\E}{\mathbb{E}}
\newcommand{\eps}{\varepsilon}
\newcommand{\dd}{\,\mathrm{d}}
\begin{document}
\section*{UnsolvedMath ID: 30003306}
\textbf{Source problem code:} catalogue entry
\subsection{Definitions and mathematical statement}
Let \(k\) denote the field appearing in the source, and let \(X\subset\mathbb{A}^{n}(k)\) be the given affine variety. No assumptions on the characteristic or algebraic closedness of \(k\), or on the dimension, irreducibility, or reducedness conventions for \(X\), are supplied. A re-embedding means a closed immersion
\[
 j\colon X\hookrightarrow\mathbb{A}^{N}_{k}
\]
for some finite integer \(N\), identifying the same variety \(X\) with its image. Here \(\mathbb{A}^{N}_{k}\) expresses the algebraic ambient space; the source writes \(\mathbb{A}^{N}(k)\) without distinguishing the variety from its set of rational points. A toric morphism, in the usual terminology, is a morphism between toric varieties whose restriction to their dense tori is a homomorphism and which is equivariant for that homomorphism. The source does not specify the toric structures or the class of toric maps allowed.

\textbf{Construction problem as supplied.} Given the affine embedding \(X\subset\mathbb{A}^{n}(k)\), without a prescribed embedded resolution, construct suitable re-embeddings \(j\colon X\hookrightarrow\mathbb{A}^{N}_{k}\) and toric maps that yield embedded local uniformizations of \(X\).

\textbf{Unresolved qualifications.} This is a construction request with an incompletely specified success condition. In the usual integral setting, local uniformization is formulated relative to a valuation ring \(V\) of the function field \(k(X)\): one seeks a suitable birational model with a center whose local ring is regular and is dominated by \(V\). Domination means that the local ring is contained in \(V\) and its maximal ideal is the contraction of the maximal ideal of \(V\). This explanation is conditional terminology, not an added integrality hypothesis or a selected interpretation of the source. The source specifies neither valuations nor centers, nor whether the construction is to work separately for each valuation or simultaneously for a collection. It also does not define the additional ambient and boundary conditions meant by \emph{embedded} local uniformization, or the domains, targets, and modification properties of the toric maps. Consequently, a fully quantified existence assertion cannot be extracted without further source clarification. The phrase ``without a prescribed embedded resolution'' states that such a resolution is not given as input; it does not assert that an embedded resolution does not exist.
\subsection{Short English statement}
Construct new affine embeddings of the given variety and suitable toric maps that provide embedded local uniformizations, without starting from a prescribed embedded resolution. The field hypotheses, valuation quantifiers, permitted toric maps, and precise embedded regularity and boundary conditions remain unspecified in the supplied record.
\subsection{Sources}
\begin{itemize}
\item Local UnsolvedMath registry, permanent ID \(30003306\), title \emph{Toric Reembeddings for Embedded Local Uniformization}, source problem code \texttt{catalogue entry}, in the Oberwolfach Reports Open Problems collection. The supplied proposed year is 2016; no individual proposer is identified.
\item Supplied report citation: \emph{Asymptotic Phenomena in Local Algebra and Singularity Theory} (2017), Oberwolfach Reports, pp.~3199--3259, DOI \texttt{10.4171/owr/2016/57}. \url{https://doi.org/10.4171/owr/2016/57}. The supplied background describes these reports as records of mathematical workshops and their open-problem discussions. No problem-specific page or passage is supplied.
\item Supplied literature source: Hauser, \emph{singularities in positive characteristic}, as labeled in the record. \url{https://arxiv.org/abs/1711.03453}. No more precise bibliographic identification or supporting passage is supplied.
\item Supplied literature source: \emph{OWR 57 (2017)}. \url{https://doi.org/10.4171/owr/2016/57}.
\item The literature assessment dated \texttt{2026-08-21} labels the problem open and reports that no general construction of the proposed re-embeddings is known that replaces the missing embedded resolution in the stated setting. This is supplied dated curation and literature triage, not an independently verified assessment; the record expressly requests independent verification of time-sensitive status and cited claims.
\item The supplied statement label is \texttt{corrected\_verified}. Its accompanying qualification says that only safe notation, typography, and TeX-delimiter cleanup was applied, preserving hypotheses, quantifiers, and subquestions. This label does not resolve the formulation gaps identified here or establish mathematical review.
\end{itemize}
\end{document}
